\section{Ubicación de la conferencia: por qué una ``receta de entrenamiento'' importa más que un truco puntual}

\subsection{Definición del problema: el razonamiento no es una capacidad única, sino un producto sistémico de la cadena de entrenamiento}
El valor central de esta conferencia no es presentar otro algoritmo nuevo, sino descomponer la ``capacidad de razonamiento'' hasta la realidad de la ingeniería: está determinada en conjunto por la distribución del preentrenamiento, los datos del posentrenamiento, el objetivo de optimización, el mecanismo de verificación y el presupuesto en tiempo de inferencia. Es decir, la capacidad de razonamiento no es el resultado directo de algún truco mágico, sino un problema de diseño de sistema de extremo a extremo.

\begin{importantbox}{Postura central del expositor}
\textit{``Open scientific research is why we are here.''} El expositor atribuye los avances de los LLM de los últimos años al ecosistema de investigación científica abierta, y enfatiza que cuando ese ecosistema se vuelve cerrado, disminuyen la reproducibilidad, la comparabilidad y la capacidad de corrección de la investigación.
\end{importantbox}

\begin{knowledgebox}{Por qué esta conferencia enfatiza las ``recipes''}
La discusión sobre la capacidad de los modelos suele reducirse a la escala de parámetros y al cómputo, pero una ``receta'' industrialmente reproducible en realidad contiene más variables implícitas:
\begin{itemize}
    \item la estrategia de muestreo y mezcla de datos;
    \item la división de las etapas de entrenamiento (preentrenamiento / posentrenamiento / tiempo de inferencia);
    \item el método de evaluación y la estrategia de descontaminación;
    \item la solución óptima de compromiso bajo restricciones de presupuesto.
\end{itemize}
\end{knowledgebox}

\begin{figure}[H]
\centering
\includegraphics[width=0.9\textwidth]{frame-01-open-science.jpg}
\caption{Video aproximadamente en 00:03:10: el expositor plantea explícitamente la ``investigación abierta'' como premisa de la metodología de la receta de entrenamiento}
\end{figure}

\subsection{De la ``publicación de modelos'' a la ``publicación científica''}
El expositor divide los modelos abiertos en tres niveles: solo pesos abiertos, código de entrenamiento abierto, y datos y proceso de evaluación abiertos. El que realmente sostiene la investigación científica es el tercer nivel, porque solo cuando se abren a la vez el procesamiento de datos, los scripts de entrenamiento y los criterios de evaluación puede la comunidad de investigación determinar de dónde proviene la mejora.

\begin{warningbox}{Abrir solo los pesos no equivale a ser reproducible}
Muchos trabajos solo publican el checkpoint, con lo cual los investigadores no pueden verificar si la mejora proviene de:
\begin{itemize}
    \item una mejora en la calidad de los datos;
    \item un ajuste fino de los hiperparámetros de entrenamiento;
    \item una modificación del optimizador o del programador;
    \item una contaminación de la evaluación o una diferencia de referencia.
\end{itemize}
Cuando estas variables no son visibles, el valor del llamado ``SOTA'' para la comunidad científica disminuye de manera notable.
\end{warningbox}

\subsection{Resumen de la sección}
Esta sección establece la línea base metodológica de toda la conferencia: la capacidad de razonamiento debe entenderse a partir de una ``receta de entrenamiento completa'', y la apertura y la reproducibilidad son condiciones necesarias para que esa receta pueda iterarse de manera acumulativa.

